% =============================================================================================
% SUPPLEMENTARY MATERIAL.  Compiles standalone: pdflatex supplementary && bibtex supplementary
% && pdflatex supplementary twice. Shares references.bib with the manuscript.
%
% The abstract of the main text promises "ten further estimates as conditional, naming the
% measurement that would qualify each". Table S1 is that promise. It did not exist until now.
%
% Sections S2-S4 were moved out of the main text's Results. They are a genuine result -- the
% argument that the barrier is now experimental rather than computational -- but a secondary one:
% they appear in neither the abstract nor the Conclusions, and at ~700 words they were the single
% largest block in the body not tied to the paper's central claim.
%
% Every number here is carried over unchanged from the verified set. Nothing was recomputed for
% this document, and nothing was added to it.
% =============================================================================================

\documentclass[preprint, 12pt]{elsarticle}

\usepackage[T1]{fontenc}
\usepackage[utf8]{inputenc}
\usepackage{amsmath, amssymb}
\usepackage{graphicx}
\usepackage{booktabs}
\usepackage[separate-uncertainty=true, detect-weight=true]{siunitx}
\usepackage[version=4]{mhchem}
\usepackage{microtype}
\usepackage[hidelinks]{hyperref}

\sisetup{group-separator = {,}, group-minimum-digits = 4}

% Number everything S1, S2, ... so nothing collides with the main text.
\renewcommand{\thesection}{S\arabic{section}}
\renewcommand{\thetable}{S\arabic{table}}
\renewcommand{\thefigure}{S\arabic{figure}}

\journal{Computational Materials Science}

\begin{document}

\begin{frontmatter}
\title{Supplementary material for\\
       \emph{Calibrating first-principles lattice thermal conductivity to experiment:
       a curated half-Heusler dataset and a validated transfer function}}
\end{frontmatter}

\section{Conditional predictions}
\label{sec:conditional}

Table~\ref{tab:conditional} lists the ten compounds that pass every structural and chemical screen
of the main text but sit in a bonding family holding no usable measured anchor. They are not
predictions. Each is the value the method \emph{would} issue if its family were validated, and the
final column names the measurement that would validate it. They are tabulated so that a reader can
see what the method is withholding and why, and so that the compounds are on record before the
measurements are made rather than after.

Nine of the ten are \ce{Ni}--\ce{Bi}, whose sole measured member is unusable for the reason set out
in Section~\ref{sec:nibi}. The tenth, \ce{HfCoBi}, sits in \ce{Co}--\ce{Bi}, which has exactly one
measured member and needs one more.

One entry carries a further caveat. \ce{Sm} can be di- or trivalent, so the valence-electron count
of \ce{SmNiBi} is not secure at 18 and its admission to the domain rests on an assumption the other
nine do not require. We flag it rather than drop it, on the same grounds as the \ce{Yb} flag in the
main text.

\begin{table}[htb]
  \centering
  \caption{Conditional estimates at \SI{300}{\kelvin}: compounds the method would predict once
  their bonding family carried two sound measurements. $\kappa^{\mathrm{BTE}}$ is the published
  transport calculation; the estimate is that value multiplied by the transfer function. Intervals
  are the conformal bands of the main text. $n_{\mathrm{src}}$ is the number of independent
  calculations available and \emph{spread} their maximum ratio, so a spread above about \num{2}
  would disqualify the row on the inter-calculation check of Section~\ref{sec:spread}.}
  \label{tab:conditional}
  \sisetup{table-format=2.2}
  \begin{tabular}{llcccccc}
    \toprule
    Compound & Family & $\kappa^{\mathrm{BTE}}$ & Estimate & \SI{50}{\percent} &
      \SI{90}{\percent} & $n_{\mathrm{src}}$ & spread \\
    & & \multicolumn{4}{c}{(\si{\watt\per\meter\per\kelvin})} & & \\
    \midrule
    \ce{PrNiBi} & \ce{Ni}--\ce{Bi} &  4.47 & 2.28 & 1.74--2.99 & 1.02--5.11  & 3 & 1.36 \\
    \ce{NdNiBi} & \ce{Ni}--\ce{Bi} &  5.25 & 2.68 & 2.04--3.51 & 1.20--6.00  & 2 & 1.00 \\
    \ce{GdNiBi} & \ce{Ni}--\ce{Bi} &  8.12 & 4.14 & 3.16--5.42 & 1.85--9.28  & 2 & 1.00 \\
    \ce{SmNiBi}$^{\dagger}$ & \ce{Ni}--\ce{Bi} &  8.49 & 4.33 & 3.31--5.67 & 1.93--9.70  & 1 & --- \\
    \ce{TbNiBi} & \ce{Ni}--\ce{Bi} &  8.57 & 4.37 & 3.34--5.73 & 1.95--9.79  & 3 & 1.31 \\
    \ce{TmNiBi} & \ce{Ni}--\ce{Bi} &  9.36 & 4.77 & 3.64--6.25 & 2.13--10.69 & 3 & 1.32 \\
    \ce{ErNiBi} & \ce{Ni}--\ce{Bi} &  9.68 & 4.94 & 3.77--6.47 & 2.20--11.06 & 3 & 1.26 \\
    \ce{LuNiBi} & \ce{Ni}--\ce{Bi} &  9.83 & 5.01 & 3.83--6.57 & 2.24--11.23 & 2 & 1.00 \\
    \ce{DyNiBi} & \ce{Ni}--\ce{Bi} & 11.78 & 6.01 & 4.59--7.87 & 2.68--13.46 & 1 & --- \\
    \midrule
    \ce{HfCoBi} & \ce{Co}--\ce{Bi} & 18.60 & 9.49 & 7.24--12.43 & 4.23--21.25 & 1 & --- \\
    \bottomrule
    \multicolumn{8}{l}{\footnotesize $^{\dagger}$ \ce{Sm} may be di- or trivalent, so
      $\mathrm{VEC}=18$ is an assumption here rather than a count.} \\
    \multicolumn{8}{l}{\footnotesize The nine \ce{Ni}--\ce{Bi} rows are unlocked together by two
      sound \ce{Ni}--\ce{Bi} measurements; \ce{HfCoBi} by one further \ce{Co}--\ce{Bi} measurement.}
  \end{tabular}
\end{table}

\section{The fourteen candidates and the measurements that unlock them}
\label{sec:fourteen}

Five laboratory measurements would take the method from five predictions to nineteen. Three
families are one measurement short: \ce{Co}--\ce{Bi}, \ce{Ni}--\ce{Pb} and \ce{Sb}--\ce{Ir} each
have exactly one measured member, and in each case the prediction for that member was correct
within a factor of two, so a second measured compound in any of them would qualify the family and
release two predictions apiece---\ce{HfCoBi} and \ce{TiCoBi}, \ce{HfNiPb} and \ce{TiNiPb},
\ce{HfSbIr} and \ce{TiSbIr}. \ce{Ni}--\ce{Bi} needs two and would release nine, which makes it by
some margin the most economical of the four.

We audited all seventeen resulting candidates against every screen that does not depend on the new
measurement. Fourteen survive. \ce{YbNiBi} fails outright: it has no cubic record, no transport
calculation, and an X site that appears nowhere in training. \ce{TiNiPb} and \ce{HoNiBi} fail the
inter-calculation check of Section~\ref{sec:spread}. Of the fourteen survivors, the ten carrying a
calculation at room temperature are tabulated above; \ce{TiCoBi}, \ce{HfNiPb}, \ce{HfSbIr} and
\ce{TiSbIr} have a transport calculation but not one at \SI{300}{\kelvin}, so no comparable estimate
can be quoted for them.

The sixteen sound candidates sit on X sites well represented in training---\ce{Ti} with \num{28}
compounds, \ce{Hf} with \num{17}, the rare earths with three to nine each---so the constraint on
them is the missing measurement and nothing else.

\section{Why \ce{Ni}--\ce{Bi} is blocked, and why reading further does not unblock it}
\label{sec:nibi}

\ce{Ni}--\ce{Bi} is blocked by an unusable reference rather than a missing one, and the distinction
matters because it is not fixed by harvesting more literature.

Its only measured member is \ce{YNiBi}, whose reported $\kappa_L$ \emph{rises} with temperature:
$\mathrm{d}\ln\kappa_L/\mathrm{d}\ln T = +\num{0.51}$ over ten points. Lattice conduction cannot
behave that way. The reason is recorded in the source itself, where the measuring authors report a
significant bipolar contribution to the thermal conductivity of \ce{YNiBi}~\cite{li2015synthesis}.
A single-parabolic-band Lorenz number carries no bipolar term, so above the bipolar onset the
electronic part is under-subtracted and the remainder booked as lattice conductivity climbs with
temperature. The defect is therefore in the electronic--lattice separation rather than in the
specimen or in our transcription---but it is disqualifying either way, and we cannot credit the
curve as validation even though our prediction for it happens to be accurate (\SI{21.7}{\percent}
error, ratio \num{0.86}).

Reading further does not close the gap. \ce{YNiBi} is the only rare-earth \ce{Ni}--\ce{Bi} half
Heusler that has ever been synthesised and measured, and of the \num{50} works citing it not one
reports a second measurement. \ce{LuNiBi}, the best studied of the candidates, has been probed by
inelastic X-ray scattering together with first-principles calculations~\cite{li2024multiphono}, but
no measured $\kappa_L$ for it exists in any source we harvested. The one \ce{Ni}--\ce{Bi} compound
that does carry an independent $\kappa_L$ measurement is \ce{ZrNiBi}, and it is disqualified on
grounds its own authors establish: stoichiometric \ce{ZrNiBi} has $\mathrm{VEC}=19$, lies above the
convex hull, and can be made phase-pure only as the \ce{Zr}-deficient \ce{Zr0.88NiBi}, so it fails
both the valence-count screen and the restriction to stoichiometric
compounds~\cite{ren2023vacancymed,liu2025stabilizat}. That its exclusion is independently motivated
by the synthesis literature is a check on the screen rather than a coincidence.

\ce{Ni}--\ce{Bi} therefore stands at zero usable members and needs two: a re-measurement of
\ce{YNiBi} below its bipolar onset together with one further \ce{Ni}--\ce{Bi} compound, or two new
compounds.

\section{Disagreement between independent calculations}
\label{sec:spread}

Auditing the candidates surfaced a check that was not being applied anywhere in the pipeline. Where
a compound has been calculated by more than one group, the published values need not agree. Across
the \num{151} half Heuslers carrying two or more independent transport calculations at
\SI{300}{\kelvin}, the median disagreement is a factor of \num{1.13}, but the tail reaches
\num{18.35}.

Two candidates fail on it. \ce{TiNiPb} carries values of \num{19.5} and
\SI{109}{\watt\per\meter\per\kelvin} from two sources, a factor of \num{5.6} apart, and
\ce{HoNiBi} differs by \num{2.37}. Neither can be predicted until the discrepancy is resolved.
Taking the first available row, as a naive pipeline would, gives a \ce{TiNiPb} prediction of
\SI{55.6}{\watt\per\meter\per\kelvin}---nearly three times the highest value ever measured for any
half Heusler.

The same check bears on the five predictions the main text does issue. \ce{ErSbPt} and \ce{HoSbPt}
rest on calculations that independent groups place \num{1.68} and \num{1.72} apart. The value we
calibrate is mid-range in both cases, but the conformal interval describes the error of the transfer
function and does not include the spread among the inputs to it, so for those two the stated
\SI{90}{\percent} band should be read as a lower bound on the true uncertainty.

\section{Which measurement is the most informative}
\label{sec:informative}

\ce{Ni}--\ce{Bi} is the most productive measurement the corpus points to. A different one is the
most informative, for a reason independent of how many predictions it releases.

\ce{Bi}--\ce{Pd} is the single family this method demonstrably fails on, under-predicting its
in-domain members by a median factor of \num{2.5}. That failure is not attributable to palladium:
\ce{Sb}--\ce{Pd} shares the \ce{Pd} site and is predicted about as well as the corpus as a
whole---\SI{29.6}{\percent} median error over its four in-domain measured members, against
\SI{33.7}{\percent} for the full blind test. Whether the \ce{Bi}--\ce{Pd} failure is a property of
bismuth on a noble-metal site, or of that particular pairing, is therefore open, and a single
\ce{Bi}--\ce{Pt} measurement would settle it.

Beyond the four families named above the economics deteriorate sharply. Of the \num{52} bonding
families never measured at all, value is concentrated in only three---\ce{Bi}--\ce{Pt},
\ce{Sn}--\ce{Au} and \ce{Pb}--\ce{Au}---each of which would need two measurements from scratch to
qualify. The other \num{46} hold one or two candidates apiece, so two measurements there would buy a
single prediction, and we do not propose them.

\bibliographystyle{elsarticle-num}
\bibliography{references}

\end{document}
